\documentclass[12pt,a4paper]{ctexart}

\usepackage[top=1.8cm,bottom=1.8cm,left=1.8cm,right=1.8cm,headheight=15pt]{geometry}
\usepackage{amsmath,amssymb}
\usepackage{enumitem}
\usepackage{xcolor}
\usepackage{fancyhdr}
\usepackage{titlesec}

\pagestyle{fancy}
\fancyhf{}
\fancyhead[L]{\textbf{曲线积分习题全解（二）}}
\fancyhead[R]{\textbf{ocr-pic-set2}}
\fancyfoot[C]{\thepage}
\renewcommand{\headrulewidth}{0.8pt}
\titleformat{\section}[hang]{\Large\bfseries}{}{0pt}{}

\newcommand{\answerbox}[1]{%
  \begin{center}\fbox{\color{red}\large\bfseries #1}\end{center}}

\begin{document}

\begin{center}{\Huge\bfseries 曲线积分习题全解（二）}\end{center}
\vspace{8pt}\hrule\vspace{8pt}

% ===================================================================
\section{第1题：正方形边界——直接法 + 格林公式}

\fbox{\parbox{\dimexpr\textwidth-2\fboxsep-2\fboxrule\relax}{%
$\displaystyle\oint_L (x^2-xy^3)dx+(y^2-2xy)dy$，
$L$ 为正方形 $O(0,0),A(2,0),B(2,2),C(0,2)$ 正向边界。}}

\subsection*{法一：直接法（逐边参数化）}

\[
\begin{aligned}
OA\;(y=0,dy=0): & \int_0^2 x^2dx = \frac{8}{3} \\[4pt]
AB\;(x=2,dx=0): & \int_0^2 (y^2-4y)dy = \Bigl[\frac{y^3}{3}-2y^2\Bigr]_0^2 = -\frac{16}{3} \\[4pt]
BC\;(y=2,dy=0): & \int_2^0 (x^2-8x)dx = -\Bigl[\frac{x^3}{3}-4x^2\Bigr]_0^2 = \frac{40}{3} \\[4pt]
CO\;(x=0,dx=0): & \int_2^0 y^2dy = -\frac{8}{3}
\end{aligned}
\]

\[
I = \frac{8}{3} - \frac{16}{3} + \frac{40}{3} - \frac{8}{3} = \frac{24}{3} = 8.
\]

\subsection*{法二：格林公式}

$P=x^2-xy^3,\;Q=y^2-2xy$，
$\dfrac{\partial Q}{\partial x}-\dfrac{\partial P}{\partial y} = (-2y) - (-3xy^2) = 3xy^2-2y$.

\[
I = \iint_{[0,2]\times[0,2]} (3xy^2-2y)\,dxdy
   = \int_0^2\Bigl[\frac{3}{2}x^2y^2-2xy\Bigr]_{x=0}^{x=2}dy
   = \int_0^2(6y^2-4y)dy
   = \Bigl[2y^3-2y^2\Bigr]_0^2 = 8.
\]

\answerbox{$I = 8$}

% ===================================================================
\section{第2题：圆周——格林公式}

\fbox{\parbox{\dimexpr\textwidth-2\fboxsep-2\fboxrule\relax}{%
$\displaystyle\oint_L xy^2\,dy - x^2y\,dx$，$L$ 为正向圆周 $x^2+y^2=a^2$。}}

$P=-x^2y,\;Q=xy^2$，
$\dfrac{\partial Q}{\partial x}=y^2,\;\dfrac{\partial P}{\partial y}=-x^2$，
$\dfrac{\partial Q}{\partial x}-\dfrac{\partial P}{\partial y}=x^2+y^2$.

\[
I = \iint_D (x^2+y^2)\,dxdy
   = \int_0^{2\pi}\int_0^a r^2\cdot r\,dr\,d\theta
   = 2\pi\cdot\frac{a^4}{4}
   = \frac{\pi a^4}{2}.
\]

\answerbox{$I = \dfrac{\pi a^4}{2}$}

% ===================================================================
\section{第3题：巧用格林公式}

\fbox{\parbox{\dimexpr\textwidth-2\fboxsep-2\fboxrule\relax}{%
$\displaystyle\oint_L \Bigl(1+\frac{1}{x^2+y^2}\Bigr)(x\,dy-y\,dx)$，
$L$ 为正向圆周 $(x-1)^2+(y-1)^2=1$。}}

$P=-y(1+\frac{1}{x^2+y^2}),\;Q=x(1+\frac{1}{x^2+y^2})$.

计算偏导数时注意 $\frac{\partial}{\partial x}\bigl(\frac{x}{x^2+y^2}\bigr)=\frac{y^2-x^2}{(x^2+y^2)^2}$：

\[
\frac{\partial Q}{\partial x}=1+\frac{1}{x^2+y^2}-\frac{2x^2}{(x^2+y^2)^2},\qquad
\frac{\partial P}{\partial y}=-1-\frac{1}{x^2+y^2}+\frac{2y^2}{(x^2+y^2)^2}.
\]

\[
\frac{\partial Q}{\partial x}-\frac{\partial P}{\partial y}
=2+\frac{2}{x^2+y^2}-\frac{2(x^2+y^2)}{(x^2+y^2)^2}=2.
\]

由于原点 $(0,0)$ 到圆心 $(1,1)$ 的距离 $\sqrt{2}>1$，原点在圆外，被积函数无奇点。

\[
I = \iint_D 2\,dxdy = 2\cdot\pi\cdot1^2 = 2\pi.
\]

\answerbox{$I = 2\pi$}

% ===================================================================
\section{第4题：用曲线积分求面积}

\fbox{\parbox{\dimexpr\textwidth-2\fboxsep-2\fboxrule\relax}{%
利用曲线积分求面积：$A=\frac{1}{2}\oint_L x\,dy-y\,dx$。\\
(1) 星形线 $x=a\cos^3t,\;y=a\sin^3t$\quad
(2) 圆 $x^2+y^2=2ax$。}}

\subsection*{(1) 星形线}

\[
dx=-3a\cos^2t\sin t\,dt,\qquad dy=3a\sin^2t\cos t\,dt.
\]

\[
\begin{aligned}
x\,dy-y\,dx &= a\cos^3t\cdot3a\sin^2t\cos t - a\sin^3t\cdot(-3a\cos^2t\sin t) \\
           &= 3a^2\cos^2t\sin^2t(\cos^2t+\sin^2t)
            = 3a^2\cos^2t\sin^2t
            = \frac{3a^2}{4}\sin^2(2t).
\end{aligned}
\]

\[
A = \frac{1}{2}\cdot\frac{3a^2}{4}\int_0^{2\pi}\sin^2(2t)\,dt
   = \frac{3a^2}{8}\cdot\pi
   = \frac{3\pi a^2}{8}.
\]

\subsection*{(2) 圆 $x^2+y^2=2ax$（即 $(x-a)^2+y^2=a^2$）}

参数化：$x=a(1+\cos\theta),\;y=a\sin\theta$，$\theta:0\to2\pi$。

\[
x\,dy-y\,dx = a(1+\cos\theta)\cdot a\cos\theta - a\sin\theta\cdot(-a\sin\theta)
            = a^2(1+\cos\theta).
\]

\[
A = \frac{1}{2}\int_0^{2\pi} a^2(1+\cos\theta)\,d\theta
   = \frac{a^2}{2}\Bigl[\theta+\sin\theta\Bigr]_0^{2\pi}
   = \pi a^2.
\]

\answerbox{$A_1 = \dfrac{3\pi a^2}{8},\qquad A_2 = \pi a^2$}

% ===================================================================
\section{第5题：路径无关性 + 求积分值}

\fbox{\parbox{\dimexpr\textwidth-2\fboxsep-2\fboxrule\relax}{%
证明 $\int_L (x^4+4xy^3)dx+(6x^2y^2-5y^4)dy$ 与路径无关，
并求 $\int_{(-2,-1)}^{(3,0)}$。}}

\subsection*{证路径无关}

$P=x^4+4xy^3,\;Q=6x^2y^2-5y^4$。

\[
\frac{\partial P}{\partial y}=12xy^2,\qquad
\frac{\partial Q}{\partial x}=12xy^2.
\]

$\dfrac{\partial P}{\partial y}=\dfrac{\partial Q}{\partial x}$，故积分与路径无关。

\subsection*{求原函数}

由 $\frac{\partial u}{\partial x}=x^4+4xy^3$，得 $u=\frac{x^5}{5}+2x^2y^3+C(y)$。

由 $\frac{\partial u}{\partial y}=6x^2y^2+C'(y)=6x^2y^2-5y^4$，得 $C'(y)=-5y^4$，$C(y)=-y^5+K$。

\[
u(x,y)=\frac{x^5}{5}+2x^2y^3-y^5+K.
\]

\subsection*{计算积分}

\[
\begin{aligned}
u(3,0) &= \frac{243}{5}, \\[4pt]
u(-2,-1) &= \frac{-32}{5}+2\cdot4\cdot(-1)-(-1)
          = -\frac{32}{5}-8+1 = -\frac{67}{5}.
\end{aligned}
\]

\[
I = u(3,0)-u(-2,-1) = \frac{243}{5}-\Bigl(-\frac{67}{5}\Bigr)
   = \frac{310}{5} = 62.
\]

\answerbox{$I = 62$}

% ===================================================================
\section{第6题：全微分验证 + 求 $u(x,y)$}

\fbox{\parbox{\dimexpr\textwidth-2\fboxsep-2\fboxrule\relax}{%
验证是全微分并求 $u(x,y)$：\\
(1) $2xy\,dx+x^2\,dy$\qquad
(2) $(2x\cos y+y^2\cos x)dx+(2y\sin x-x^2\sin y)dy$。}}

\subsection*{(1)}

$P=2xy,\;Q=x^2$。
$\frac{\partial P}{\partial y}=2x=\frac{\partial Q}{\partial x}$ ✓

由 $\frac{\partial u}{\partial x}=2xy$，得 $u=x^2y+C(y)$。
由 $\frac{\partial u}{\partial y}=x^2+C'(y)=x^2$，得 $C(y)=K$。

\[
\boxed{u(x,y)=x^2y+K}
\]

\subsection*{(2)}

$P=2x\cos y+y^2\cos x,\;Q=2y\sin x-x^2\sin y$。

\[
\frac{\partial P}{\partial y}=-2x\sin y+2y\cos x,\qquad
\frac{\partial Q}{\partial x}=2y\cos x-2x\sin y.
\]

$\frac{\partial P}{\partial y}=\frac{\partial Q}{\partial x}$ ✓

由 $\frac{\partial u}{\partial x}=2x\cos y+y^2\cos x$，得 $u=x^2\cos y+y^2\sin x+C(y)$。
由 $\frac{\partial u}{\partial y}=-x^2\sin y+2y\sin x+C'(y)=2y\sin x-x^2\sin y$，得 $C(y)=K$。

\[
\boxed{u(x,y)=x^2\cos y+y^2\sin x+K}
\]

\answerbox{全微分验证通过，原函数如上。}

% ===================================================================
\section{B类 第1题：路径无关 + 求积分}

\fbox{\parbox{\dimexpr\textwidth-2\fboxsep-2\fboxrule\relax}{%
$\displaystyle\int_L (y^2e^x+3x^2+2xy+y^2)dx+(2ye^x+x^2+2xy-3y^2)dy$，
$L$ 为上半圆 $y=\sqrt{1-x^2}$ 从 $A(1,0)$ 到 $B(0,1)$。}}

\subsection*{验证路径无关}

$\frac{\partial P}{\partial y}=2ye^x+2x+2y,\qquad
\frac{\partial Q}{\partial x}=2ye^x+2x+2y$。相等 ✓

\subsection*{求原函数}

由 $\frac{\partial u}{\partial x}=y^2e^x+3x^2+2xy+y^2$，
$u=y^2e^x+x^3+x^2y+xy^2+C(y)$。

$\frac{\partial u}{\partial y}=2ye^x+x^2+2xy+C'(y)=2ye^x+x^2+2xy-3y^2$，
$C'(y)=-3y^2,\;C(y)=-y^3+K$。

\[
u=y^2e^x+x^3+x^2y+xy^2-y^3+K.
\]

\subsection*{计算}

\[
u(0,1)=1\cdot1+0+0+0-1=0,\qquad
u(1,0)=0+1+0+0-0=1.
\]

\[
I = u(0,1)-u(1,0)=0-1=-1.
\]

\answerbox{$I = -1$}

% ===================================================================
\section{B类 第2题：经典积分——分类讨论}

\fbox{\parbox{\dimexpr\textwidth-2\fboxsep-2\fboxrule\relax}{%
$\displaystyle I=\oint_L \frac{x\,dy-y\,dx}{x^2+y^2}$，
$L$ 为正向圆周 $x^2+(y-1)^2=R^2\;(R\neq1)$。}}

\subsection*{理论分析}

\[
P=-\frac{y}{x^2+y^2},\qquad Q=\frac{x}{x^2+y^2}.
\]

当 $(x,y)\neq(0,0)$ 时：
\[
\frac{\partial P}{\partial y}=-\frac{x^2-y^2}{(x^2+y^2)^2}=\frac{\partial Q}{\partial x}.
\]

\textbf{关键是原点 $(0,0)$ 是否在曲线内：}
圆心 $(0,1)$ 到原点的距离为 $1$。

\begin{itemize}[leftmargin=*]
  \item 若 $0<R<1$：曲线不包围原点。由格林公式，$I=0$。
  \item 若 $R>1$：曲线包围原点。绕原点的闭路积分为 $2\pi$（在小圆 $x^2+y^2=\varepsilon^2$ 上直接计算可得）。
\end{itemize}

\subsection*{$R>1$ 情形的计算}

在圆周 $x^2+y^2=a^2$ 上取正向：
\[
x=a\cos t,\;y=a\sin t,\;t:0\to2\pi.
\]
\[
\frac{x\,dy-y\,dx}{x^2+y^2}
   =\frac{a\cos t\cdot a\cos t - a\sin t\cdot(-a\sin t)}{a^2}
   = \cos^2t+\sin^2t = 1.
\]
\[
\oint_{x^2+y^2=a^2} \frac{x\,dy-y\,dx}{x^2+y^2}
   = \int_0^{2\pi} 1\,dt = 2\pi.
\]

此结果与半径 $a$ 无关，任何包围原点的正向闭曲线的积分值都是 $2\pi$。

\answerbox{$I = \begin{cases} 0, & 0<R<1 \text{（不包围原点）}\\ 2\pi, & R>1 \text{（包围原点）} \end{cases}$}

\end{document}
